\documentclass[11pt,a4paper]{article}

% ---- Packages ----
\usepackage[utf8]{inputenc}
\usepackage[T1]{fontenc}
\usepackage{amsmath,amssymb,amsthm}
\usepackage{mathtools}
\usepackage[margin=1in,headheight=14pt]{geometry}
\usepackage{hyperref}
\usepackage{enumitem}
\usepackage{xcolor}
\usepackage{tcolorbox}
\usepackage{fancyhdr}
\usepackage{booktabs}

% ---- Hyperref ----
\hypersetup{
    colorlinks=true,
    linkcolor=red,
    citecolor=blue,
    urlcolor=blue
}

% ---- Header ----
\pagestyle{fancy}
\fancyhf{}
\fancyhead[L]{\textit{Lesson 6: Metric and Normed Spaces}}
\fancyhead[R]{\thepage}
\renewcommand{\headrulewidth}{0.4pt}

% ---- Theorem Environments ----
\theoremstyle{definition}
\newtheorem{definition}{Definition}[section]
\newtheorem{example}{Example}[section]
\newtheorem{exercise}{Exercise}[section]

\theoremstyle{plain}
\newtheorem{theorem}{Theorem}[section]
\newtheorem{lemma}[theorem]{Lemma}
\newtheorem{proposition}[theorem]{Proposition}
\newtheorem{corollary}[theorem]{Corollary}

\theoremstyle{remark}
\newtheorem{remark}{Remark}[section]

% ---- Colored Boxes ----
\tcbuselibrary{theorems,skins,breakable}

\newtcolorbox{keyresult}[1][]{
    colback=green!5!white,
    colframe=green!50!black,
    fonttitle=\bfseries,
    title=#1,
    breakable
}

\newtcolorbox{rlconnection}[1][]{
    colback=blue!5!white,
    colframe=blue!50!black,
    fonttitle=\bfseries,
    title=#1,
    breakable
}

\newtcolorbox{warning}[1][]{
    colback=red!5!white,
    colframe=red!50!black,
    fonttitle=\bfseries,
    title=#1,
    breakable
}

% ---- Operators ----
\newcommand{\R}{\mathbb{R}}
\newcommand{\N}{\mathbb{N}}
\newcommand{\Q}{\mathbb{Q}}
\newcommand{\Z}{\mathbb{Z}}
\newcommand{\C}{\mathbb{C}}
\newcommand{\Prob}{\mathbb{P}}
\newcommand{\E}{\mathbb{E}}
\newcommand{\indicator}{\mathbf{1}}
\DeclareMathOperator{\diam}{diam}
\DeclareMathOperator{\Int}{int}
\DeclareMathOperator{\cl}{cl}
\DeclareMathOperator{\Span}{span}
\DeclareMathOperator{\id}{id}

% ---- Title ----
\title{Lesson 15: Metric Spaces}
\author{Curriculum: A Rigorous Learning Path to Multi-Agent Deep RL}
\date{}

\begin{document}
\maketitle
\tableofcontents
\newpage

\setcounter{section}{0}

\section{Why Functional Analysis?}
\label{sec:intro}

Reinforcement learning algorithms operate on \emph{function spaces}. A value
function $V \colon \mathcal{S} \to \R$ assigns a real number to every state; a
Q-function $Q \colon \mathcal{S} \times \mathcal{A} \to \R$ assigns a value to
every state-action pair. When we write the Bellman optimality equation
\begin{equation}
\label{eq:bellman_intro}
V^*(s) = \max_{a \in \mathcal{A}} \left[ r(s,a) + \gamma \sum_{s'} p(s' \mid s,a)\, V^*(s') \right],
\end{equation}
we are asserting that $V^*$ is a \emph{fixed point} of an operator $T$ acting on
a space of functions. The fundamental questions are:

\begin{enumerate}[nosep]
    \item Does a fixed point exist?
    \item Is it unique?
    \item Does the iteration $V_{n+1} = T V_n$ converge to $V^*$?
    \item How fast does it converge?
\end{enumerate}

Answering these questions requires the language of functional analysis:

\begin{itemize}[nosep]
    \item \textbf{Metric spaces} (Section~\ref{sec:metric}) provide the notion of
    distance needed to say ``$V_n$ converges to $V^*$.''
    \item \textbf{Normed spaces} (Section~\ref{sec:normed}) give us a notion of
    ``size'' for functions and a way to measure error $\|V_n - V^*\|$.
    \item \textbf{Completeness} (Section~\ref{sec:banach}) ensures that Cauchy
    sequences converge, which is essential for the Banach fixed-point theorem.
    \item \textbf{Bounded linear operators} (Section~\ref{sec:operators}) formalize
    the maps between function spaces that arise in approximation theory.
    \item \textbf{Compactness} (Section~\ref{sec:compactness}) is the key to
    existence results and understanding the behavior of infinite-dimensional spaces.
    \item \textbf{The Baire category theorem} (Section~\ref{sec:baire}) underpins
    the three pillars of linear functional analysis: the Uniform Boundedness
    Principle, the Open Mapping Theorem, and the Closed Graph Theorem.
\end{itemize}

\begin{rlconnection}[The Bellman Operator as Motivation]
The Bellman optimality operator $T$ defined by the right-hand side of
\eqref{eq:bellman_intro} acts on the Banach space $(b(\mathcal{S}), \|\cdot\|_\infty)$
of bounded real-valued functions on the state space, equipped with the supremum
norm. The contraction mapping theorem (Banach fixed-point theorem, which we prove
in the next lesson) guarantees that $T$ has a unique fixed point and that
$\|T^n V_0 - V^*\|_\infty \le \gamma^n \|V_0 - V^*\|_\infty$. Understanding
\emph{why} this works requires every concept in this lesson: metric, norm,
completeness, and the operator framework.
\end{rlconnection}

%% ============================================================
%% SECTION 2: METRIC SPACES
%% ============================================================

\section{Metric Spaces}
\label{sec:metric}

\subsection{Definition and Examples}
\label{subsec:metric_def}

\begin{definition}[Metric Space]
\label{def:metric}
A \emph{metric} on a set $X$ is a function $d \colon X \times X \to [0,\infty)$
satisfying, for all $x, y, z \in X$:
\begin{enumerate}[nosep,label=(M\arabic*)]
    \item \label{M1} $d(x,y) = 0$ if and only if $x = y$.
    \hfill(identity of indiscernibles)
    \item \label{M2} $d(x,y) = d(y,x)$.
    \hfill(symmetry)
    \item \label{M3} $d(x,z) \le d(x,y) + d(y,z)$.
    \hfill(triangle inequality)
\end{enumerate}
The pair $(X,d)$ is called a \emph{metric space}. Elements of $X$ are called
\emph{points}.
\end{definition}

\begin{remark}
Non-negativity $d(x,y) \ge 0$ follows from the axioms: by \ref{M3} and \ref{M1},
$0 = d(x,x) \le d(x,y) + d(y,x) = 2d(x,y)$ by \ref{M2}, so $d(x,y) \ge 0$.
\end{remark}

\begin{example}[Euclidean Metric]
\label{ex:euclidean}
$X = \R^n$ with $d(x,y) = \bigl(\sum_{i=1}^n (x_i - y_i)^2\bigr)^{1/2}$. The
triangle inequality follows from the Cauchy--Schwarz inequality.
\end{example}

\begin{example}[$p$-Metrics on $\R^n$]
\label{ex:p_metric}
For $1 \le p \le \infty$, define $d_p(x,y) = \|x - y\|_p$ where
\[
\|x\|_p =
\begin{cases}
\bigl(\sum_{i=1}^n |x_i|^p\bigr)^{1/p} & \text{if } 1 \le p < \infty,\\
\max_{1 \le i \le n} |x_i| & \text{if } p = \infty.
\end{cases}
\]
Each $d_p$ is a metric. The triangle inequality for $1 < p < \infty$ is
Minkowski's inequality (proved in Proposition~\ref{prop:minkowski}).
\end{example}

\begin{example}[Discrete Metric]
\label{ex:discrete}
On any set $X$, define $d(x,y) = 0$ if $x = y$ and $d(x,y) = 1$ if $x \ne y$.
This is a metric. Every subset of $X$ is open (and closed), so $X$ with the
discrete metric is a \emph{discrete topological space}.
\end{example}

\begin{example}[Function Space $C[a,b]$]
\label{ex:Cab}
Let $C[a,b]$ denote the set of continuous real-valued functions on $[a,b]$. Define
\[
d_\infty(f,g) = \max_{t \in [a,b]} |f(t) - g(t)|, \qquad
d_1(f,g) = \int_a^b |f(t) - g(t)|\, dt.
\]
Both are metrics on $C[a,b]$. The metric $d_\infty$ measures uniform distance;
$d_1$ measures average distance.
\end{example}

\begin{example}[Sequence Spaces $\ell^p$]
\label{ex:ell_p}
For $1 \le p < \infty$, define $\ell^p = \{(x_n)_{n=1}^\infty \subset \R :
\sum_{n=1}^\infty |x_n|^p < \infty\}$ with metric
$d(x,y) = \bigl(\sum_{n=1}^\infty |x_n - y_n|^p\bigr)^{1/p}$.
For $p = \infty$, $\ell^\infty = \{(x_n) : \sup_n |x_n| < \infty\}$ with
$d(x,y) = \sup_n |x_n - y_n|$.
\end{example}

\begin{rlconnection}[State Spaces as Metric Spaces]
In RL with a finite state space $\mathcal{S} = \{s_1, \ldots, s_n\}$, a value
function $V \colon \mathcal{S} \to \R$ is simply a vector in $\R^n$. The space of
all value functions is $(\R^n, \|\cdot\|_\infty)$, which is a metric space. For
infinite or continuous state spaces, value functions live in function spaces like
$C(\mathcal{S})$ or $\ell^\infty(\mathcal{S})$, and the theory of this section
becomes essential.
\end{rlconnection}

\subsection{Topology of Metric Spaces}
\label{subsec:topology}

\begin{definition}[Open Ball]
\label{def:open_ball}
In a metric space $(X,d)$, the \emph{open ball} of radius $r > 0$ centered at
$x \in X$ is
\[
B(x,r) = \{y \in X : d(x,y) < r\}.
\]
The \emph{closed ball} is $\overline{B}(x,r) = \{y \in X : d(x,y) \le r\}$.
\end{definition}

\begin{definition}[Open and Closed Sets]
\label{def:open_closed}
Let $(X,d)$ be a metric space.
\begin{enumerate}[nosep,label=(\roman*)]
    \item A set $G \subseteq X$ is \emph{open} if for every $x \in G$ there exists
    $r > 0$ such that $B(x,r) \subseteq G$.
    \item A set $F \subseteq X$ is \emph{closed} if its complement $X \setminus F$
    is open.
\end{enumerate}
\end{definition}

\begin{proposition}[Properties of Open and Closed Sets]
\label{prop:open_closed}
In a metric space $(X,d)$:
\begin{enumerate}[nosep,label=(\alph*)]
    \item $\emptyset$ and $X$ are both open and closed.
    \item Any union of open sets is open.
    \item Any finite intersection of open sets is open.
    \item Any intersection of closed sets is closed.
    \item Any finite union of closed sets is closed.
    \item Every open ball $B(x,r)$ is open.
\end{enumerate}
\end{proposition}

\begin{proof}
\textbf{(a)} $\emptyset$ is open vacuously. $X$ is open because $B(x,r) \subseteq X$
for all $x$ and $r > 0$.

\textbf{(b)} Let $\{G_\alpha\}_{\alpha \in I}$ be a family of open sets and let
$x \in \bigcup_\alpha G_\alpha$. Then $x \in G_\beta$ for some $\beta$, so there
exists $r > 0$ with $B(x,r) \subseteq G_\beta \subseteq \bigcup_\alpha G_\alpha$.

\textbf{(c)} Let $G_1, \ldots, G_n$ be open and let $x \in \bigcap_{i=1}^n G_i$.
For each $i$, there exists $r_i > 0$ with $B(x,r_i) \subseteq G_i$. Set
$r = \min(r_1, \ldots, r_n) > 0$. Then $B(x,r) \subseteq \bigcap_{i=1}^n G_i$.

\textbf{(d)--(e)} Follow from (b)--(c) by taking complements (De Morgan's laws).

\textbf{(f)} Let $y \in B(x,r)$ and set $\varepsilon = r - d(x,y) > 0$. For any
$z \in B(y,\varepsilon)$, we have $d(x,z) \le d(x,y) + d(y,z) < d(x,y) +
\varepsilon = r$, so $z \in B(x,r)$. Thus $B(y,\varepsilon) \subseteq B(x,r)$.
\end{proof}

\begin{definition}[Interior, Closure, Boundary]
\label{def:int_cl_bd}
Let $A \subseteq X$ in a metric space $(X,d)$.
\begin{enumerate}[nosep,label=(\roman*)]
    \item The \emph{interior} of $A$ is $\Int(A) = \bigcup \{G \subseteq A : G \text{ is open}\}$.
    \item The \emph{closure} of $A$ is $\cl(A) = \bigcap \{F \supseteq A : F \text{ is closed}\}$.
    \item The \emph{boundary} of $A$ is $\partial A = \cl(A) \setminus \Int(A)$.
    \item $A$ is \emph{dense} in $X$ if $\cl(A) = X$.
\end{enumerate}
\end{definition}

\begin{proposition}[Closure Characterization]
\label{prop:closure_char}
$x \in \cl(A)$ if and only if for every $r > 0$, $B(x,r) \cap A \ne \emptyset$.
Equivalently, $x \in \cl(A)$ if and only if there exists a sequence $(x_n)$ in $A$
with $x_n \to x$.
\end{proposition}

\begin{proof}
($\Rightarrow$) Suppose $x \in \cl(A)$ and there exists $r > 0$ with
$B(x,r) \cap A = \emptyset$. Then $A \subseteq X \setminus B(x,r)$, which is
closed. So $\cl(A) \subseteq X \setminus B(x,r)$, contradicting $x \in \cl(A)$.

($\Leftarrow$) Suppose every ball around $x$ meets $A$, and let $F$ be any closed
set containing $A$. If $x \notin F$, then $x \in X \setminus F$, which is open,
so there exists $r > 0$ with $B(x,r) \subseteq X \setminus F$. But then
$B(x,r) \cap A \subseteq B(x,r) \cap F = \emptyset$, contradicting our assumption.
So $x \in F$ for every closed $F \supseteq A$, hence $x \in \cl(A)$.

For the sequential characterization: if $B(x,1/n) \cap A \ne \emptyset$ for each
$n$, pick $x_n \in B(x,1/n) \cap A$; then $d(x_n,x) < 1/n \to 0$. Conversely, if
$x_n \to x$ with each $x_n \in A$, then every ball $B(x,r)$ contains $x_n$ for
large $n$, so $B(x,r) \cap A \ne \emptyset$.
\end{proof}

\subsection{Convergence and Completeness}
\label{subsec:convergence}

\begin{definition}[Convergence]
\label{def:convergence}
A sequence $(x_n)$ in a metric space $(X,d)$ \emph{converges} to $x \in X$ if
$d(x_n, x) \to 0$ as $n \to \infty$, i.e., for every $\varepsilon > 0$ there
exists $N \in \N$ such that $d(x_n, x) < \varepsilon$ for all $n \ge N$. We write
$x_n \to x$ or $\lim_{n \to \infty} x_n = x$.
\end{definition}

\begin{proposition}[Uniqueness of Limits]
\label{prop:unique_limit}
In a metric space, limits are unique: if $x_n \to x$ and $x_n \to y$, then $x = y$.
\end{proposition}

\begin{proof}
$0 \le d(x,y) \le d(x,x_n) + d(x_n,y) \to 0 + 0 = 0$, so $d(x,y) = 0$, hence $x = y$.
\end{proof}

\begin{definition}[Cauchy Sequence]
\label{def:cauchy}
A sequence $(x_n)$ in $(X,d)$ is a \emph{Cauchy sequence} if for every
$\varepsilon > 0$ there exists $N \in \N$ such that $d(x_m, x_n) < \varepsilon$
for all $m, n \ge N$.
\end{definition}

\begin{proposition}
\label{prop:convergent_is_cauchy}
Every convergent sequence is Cauchy.
\end{proposition}

\begin{proof}
If $x_n \to x$, then for $\varepsilon > 0$ there exists $N$ such that
$d(x_n, x) < \varepsilon/2$ for $n \ge N$. For $m, n \ge N$:
$d(x_m, x_n) \le d(x_m, x) + d(x, x_n) < \varepsilon/2 + \varepsilon/2 = \varepsilon$.
\end{proof}

\begin{warning}[Cauchy Does Not Imply Convergent]
The converse fails in general. In $\Q$ with the usual metric, the sequence of
rational approximations to $\sqrt{2}$ (e.g., defined by $x_1 = 1$,
$x_{n+1} = (x_n + 2/x_n)/2$) is Cauchy but does not converge in $\Q$.
The ``missing limit'' is $\sqrt{2} \notin \Q$.
\end{warning}

\begin{definition}[Complete Metric Space]
\label{def:complete}
A metric space $(X,d)$ is \emph{complete} if every Cauchy sequence in $X$
converges to a limit in $X$.
\end{definition}

\begin{theorem}[Completeness of $\R$]
\label{thm:R_complete}
$(\R, |\cdot|)$ is a complete metric space.
\end{theorem}

\begin{proof}
Let $(x_n)$ be a Cauchy sequence in $\R$. First, $(x_n)$ is bounded: there exists
$N$ such that $|x_m - x_n| < 1$ for all $m, n \ge N$, so
$|x_n| \le |x_N| + 1$ for $n \ge N$, giving
$|x_n| \le \max(|x_1|, \ldots, |x_{N-1}|, |x_N| + 1)$ for all $n$.

By the Bolzano--Weierstrass theorem (every bounded sequence in $\R$ has a
convergent subsequence), there is a subsequence $x_{n_k} \to x$ for some $x \in \R$.

We show $x_n \to x$. Given $\varepsilon > 0$, choose $N_1$ such that
$|x_m - x_n| < \varepsilon/2$ for $m, n \ge N_1$, and $N_2$ such that
$|x_{n_k} - x| < \varepsilon/2$ for $k \ge N_2$. For $n \ge N_1$, pick $k$ with
$n_k \ge \max(N_1, n_{N_2})$. Then
$|x_n - x| \le |x_n - x_{n_k}| + |x_{n_k} - x| < \varepsilon/2 + \varepsilon/2 = \varepsilon$.
\end{proof}

\begin{theorem}[Completeness of $\R^n$]
\label{thm:Rn_complete}
$(\R^n, d_p)$ is complete for every $1 \le p \le \infty$.
\end{theorem}

\begin{proof}
Let $(x^{(k)})_{k=1}^\infty$ be a Cauchy sequence in $(\R^n, d_p)$, where
$x^{(k)} = (x_1^{(k)}, \ldots, x_n^{(k)})$. For each coordinate $i$,
$|x_i^{(k)} - x_i^{(m)}| \le \|x^{(k)} - x^{(m)}\|_p$ (since the $\ell^p$ norm
dominates each coordinate), so $(x_i^{(k)})_{k=1}^\infty$ is Cauchy in $\R$.
By Theorem~\ref{thm:R_complete}, $x_i^{(k)} \to x_i$ for some $x_i \in \R$.
Set $x = (x_1, \ldots, x_n)$. Then
$\|x^{(k)} - x\|_p = \bigl(\sum_{i=1}^n |x_i^{(k)} - x_i|^p\bigr)^{1/p} \to 0$
as $k \to \infty$ (each term goes to $0$; this is a finite sum). So $x^{(k)} \to x$.
\end{proof}

\begin{definition}[Dense Subset, Separable Space]
\label{def:separable}
A metric space $(X,d)$ is \emph{separable} if it contains a countable dense subset,
i.e., there exists a countable set $D \subseteq X$ with $\cl(D) = X$.
\end{definition}

\begin{example}
$\R^n$ is separable: $\Q^n$ is countable and dense. The space $C[a,b]$ with the
sup metric is separable: by the Weierstrass approximation theorem, polynomials with
rational coefficients form a countable dense subset.
\end{example}

\begin{theorem}[Completion of a Metric Space]
\label{thm:completion}
Every metric space $(X,d)$ has a \emph{completion}: there exists a complete metric
space $(\hat{X}, \hat{d})$ and an isometry $\iota \colon X \to \hat{X}$ (i.e.,
$\hat{d}(\iota(x), \iota(y)) = d(x,y)$ for all $x, y \in X$) such that
$\iota(X)$ is dense in $\hat{X}$. The completion is unique up to isometry.
\end{theorem}

\begin{proof}
\textbf{Construction.}
Let $\mathcal{C}$ be the set of all Cauchy sequences in $X$. Define an equivalence
relation: $(x_n) \sim (y_n)$ if $d(x_n, y_n) \to 0$. Let
$\hat{X} = \mathcal{C}/{\sim}$ be the set of equivalence classes, and write
$[(x_n)]$ for the class of $(x_n)$.

\textbf{The metric is well defined.}
For $(x_n), (y_n) \in \mathcal{C}$, the sequence $(d(x_n, y_n))$ is Cauchy in $\R$
(since $|d(x_m,y_m) - d(x_n,y_n)| \le d(x_m,x_n) + d(y_m,y_n)$ by the triangle
inequality), hence convergent. Define
$\hat{d}([(x_n)], [(y_n)]) = \lim_{n \to \infty} d(x_n, y_n)$.
One checks that this is independent of representatives (if $(x_n) \sim (x_n')$
and $(y_n) \sim (y_n')$, then $|d(x_n,y_n) - d(x_n',y_n')| \le d(x_n,x_n') +
d(y_n,y_n') \to 0$) and satisfies the metric axioms.

\textbf{Isometric embedding.}
Define $\iota(x) = [(x, x, x, \ldots)]$ (the class of the constant sequence).
Then $\hat{d}(\iota(x), \iota(y)) = \lim_n d(x,y) = d(x,y)$, so $\iota$ is an
isometry.

\textbf{Density of $\iota(X)$.}
Given $[(x_n)] \in \hat{X}$ and $\varepsilon > 0$, choose $N$ with
$d(x_m, x_n) < \varepsilon/2$ for $m, n \ge N$. Then
$\hat{d}([(x_n)], \iota(x_N)) = \lim_n d(x_n, x_N) \le \varepsilon/2 < \varepsilon$.

\textbf{Completeness of $\hat{X}$.}
Let $(p_k)$ be Cauchy in $\hat{X}$. Since $\iota(X)$ is dense, for each $k$ choose
$z_k \in X$ with $\hat{d}(p_k, \iota(z_k)) < 1/k$. We claim $(z_k)$ is Cauchy in
$X$: indeed, $d(z_k, z_m) = \hat{d}(\iota(z_k), \iota(z_m)) \le
\hat{d}(\iota(z_k), p_k) + \hat{d}(p_k, p_m) + \hat{d}(p_m, \iota(z_m)) <
1/k + \hat{d}(p_k, p_m) + 1/m \to 0$. Let $p = [(z_k)] \in \hat{X}$. Then
$\hat{d}(p_k, p) \le \hat{d}(p_k, \iota(z_k)) + \hat{d}(\iota(z_k), p) <
1/k + \hat{d}(\iota(z_k), p) \to 0$ since
$\hat{d}(\iota(z_k), p) = \lim_n d(z_k, z_n) \to 0$ as $k \to \infty$.

\textbf{Uniqueness.}
If $(\hat{X}_1, \hat{d}_1, \iota_1)$ and $(\hat{X}_2, \hat{d}_2, \iota_2)$ are
two completions, define $\varphi \colon \iota_1(X) \to \iota_2(X)$ by
$\varphi(\iota_1(x)) = \iota_2(x)$. This is an isometry, hence uniformly
continuous, and extends uniquely to an isometry $\hat{X}_1 \to \hat{X}_2$ by
density and completeness.
\end{proof}

\begin{rlconnection}[Completion and Function Approximation]
In RL with function approximation, we often work with a parametric class of
functions (e.g., neural networks) that is not complete under the sup norm. The
completion theorem guarantees that this class can be embedded into a complete
space. Understanding which space we ``complete to'' is crucial for analyzing
whether approximate value iteration converges.
\end{rlconnection}

%% ============================================================
%% SECTION 3: NORMED SPACES
%% ============================================================


\end{document}